% Every source of time on this part, what it is good for, and what it costs.
% Middle and right columns share one pitch of 2.4 cm so that a four-line block
% cannot touch its neighbour; the two cross-relations are routed outside the
% columns rather than between them.
\begin{tikzpicture}[font=\sffamily\small, x=1cm, y=1cm]

% ---------------- oscillators
\node[hw, text width=26mm] (hse) at (-6.2,3.9) {8 MHz from the probe\\{\scriptsize bypass, not a crystal}};
\node[hw, text width=26mm] (lse) at (-6.2,1.1) {32.768 kHz crystal\\{\scriptsize fitted, confirmed}};
\begin{scope}[on background layer]
\node[layer, fit=(hse)(lse), inner sep=6pt,
      label={[layerlabel, anchor=south west]north west:two independent oscillators}] {};
\end{scope}

% ---------------- the counters
\node[fw, text width=32mm] (dwt)  at (-0.4, 5.6) {cycle counter\\{\scriptsize 3.57 ns, wraps in 15.3 s}};
\node[fw, text width=32mm] (tim)  at (-0.4, 3.2) {TIM2, 32 bit at 1 MHz\\{\scriptsize wraps in 71.6 min, extended to 64 bit in software}};
\node[fw, text width=32mm] (rtc)  at (-0.4, 0.8) {real-time clock\\{\scriptsize seconds and a subsecond register, survives reset}};
\node[fw, text width=32mm] (tscv) at (-0.4,-1.6) {bus timestamp unit\\{\scriptsize captured at the start-of-frame bit, 16 bit}};

% ---------------- what each is for
\node[user, text width=30mm] (exec)  at (5.4, 5.6) {execution time\\{\scriptsize inside one period only}};
\node[user, text width=30mm] (mono)  at (5.4, 3.2) {the monotonic clock\\{\scriptsize intervals, ordering, every timestamp in a frame}};
\node[user, text width=30mm] (wall)  at (5.4, 0.8) {the wall clock\\{\scriptsize log lines and dates, never an interval}};
\node[user, text width=30mm] (agree) at (5.4,-1.6) {agreement with other nodes\\{\scriptsize chapter 12}};

% ---------------- what this bench cannot do
\node[absentblk, text width=30mm] (ptp) at (5.4,-4.1)
  {precision time protocol\\{\scriptsize needs an Ethernet controller with hardware timestamping}};
\node[absentblk, text width=30mm] (ttcan) at (5.4,-6.6)
  {time-triggered bus\\{\scriptsize lives in dedicated transceivers, absent from the flexible-data format}};
\begin{scope}[on background layer]
\node[layer, fit=(ptp)(ttcan), inner sep=6pt,
      label={[layerlabel, anchor=north west]south west:not available on this part}] {};
\end{scope}

% ---------------- edges
\draw[flow] (hse) -- (dwt);
\draw[flow] (hse) -- (tim);
\draw[flow] (hse) -- (tscv);
\draw[flow] (lse) -- (rtc);
\draw[flow] (dwt)  -- (exec);
\draw[flow] (tim)  -- (mono);
\draw[flow] (rtc)  -- (wall);
\draw[flow] (tscv) -- (agree);

% the two cross-relations, routed clear of both columns
\draw[hiz, <->] (tim.east) .. controls +(1.4,0) and +(1.4,0) .. (rtc.east);
\node[font=\sffamily\scriptsize, fill=white, inner sep=1.5pt, anchor=west]
  at (2.5,2.0) {drift, measured in ppm};

\draw[hiz, <->] (tim.west) .. controls +(-1.6,0) and +(-1.6,0) .. (tscv.west);
\node[font=\sffamily\scriptsize, fill=white, inner sep=1.5pt, anchor=east]
  at (-2.6,-0.6) {different clock domains};

\node[umlnote, text width=40mm, anchor=north west] at (-7.0,-2.6)
  {The node uses three sources for three jobs rather than making one do
   everything. The rule is four lines in the repository and every later chapter
   is held to it.};
\node[umlnote, text width=40mm, anchor=north west] at (-7.0,-5.0)
  {The capture in the bus controller is the useful one and it is free: the
   counter value is taken at the frame itself, so interrupt latency moves when
   the node learns about a frame, not when the frame is recorded. Nothing in
   software can match that, and it is unused until chapter 12.};
\end{tikzpicture}
